\newglossaryentry{id1}{
    name={Аномалия},
    description={отклонение наблюдаемого поведения объекта, процесса или системы от ожидаемого нормального состояния.}
}

\newglossaryentry{id2}{
    name={Аудит-лог (audit log)},
    description={журнал событий безопасности, содержащий сведения о действиях пользователей, сервисных аккаунтов и компонентов системы при обращении к API Kubernetes.}
}

\newglossaryentry{id3}{
    name={Автокодировщик (autoencoder)},
    description={нейронная сеть, обучаемая восстанавливать входные данные через компактное латентное представление, применяемая для извлечения признаков и обнаружения аномалий.}
}

\newglossaryentry{id4}{
    name={Валидационная выборка},
    description={часть набора данных, используемая для контроля процесса обучения модели и настройки её гиперпараметров.}
}

\newglossaryentry{id5}{
    name={Выборка (dataset)},
    description={подмножество данных, используемое для обучения, валидации или тестирования модели машинного обучения.}
}

\newglossaryentry{id6}{
    name={Кластер Kubernetes},
    description={совокупность узлов и управляющих компонентов, обеспечивающих развертывание, масштабирование и выполнение контейнеризированных приложений.}
}

\newglossaryentry{id7}{
    name={Контейнер (container)},
    description={изолированная среда выполнения приложения, включающая исполняемый код, библиотеки и зависимости и использующая ядро хост-системы.}
}

\newglossaryentry{id8}{
    name={Контейнеризация (containerization)},
    description={технология упаковки и запуска приложений в контейнерах, обеспечивающая переносимость и независимость от инфраструктуры.}
}

\newglossaryentry{id9}{
    name={Контроллер (controller)},
    description={компонент плоскости управления Kubernetes, отслеживающий текущее состояние объектов кластера и приводящий его к заданному состоянию.}
}

\newglossaryentry{id10}{
    name={Машинное обучение (machine learning)},
    description={раздел искусственного интеллекта, в котором модели обучаются на данных для выявления закономерностей, прогнозирования и классификации.}
}

\newglossaryentry{id11}{
    name={Нейронная сеть (neural network)},
    description={метод машинного обучения, основанный на системе взаимосвязанных искусственных нейронов и применяемый для аппроксимации сложных зависимостей в данных.}
}

\newglossaryentry{id12}{
    name={Обучающая выборка},
    description={часть набора данных, используемая для настройки параметров модели в процессе обучения.}
}

\newglossaryentry{id13}{
    name={Оркестрация},
    description={процесс автоматизации развертывания, масштабирования и управления контейнеризованными приложениями с учётом зависимостей и отказоустойчивости.}
}

\newglossaryentry{id14}{
    name={Плоскость управления (control plane)},
    description={набор компонентов Kubernetes, отвечающих за управление состоянием кластера, планирование задач и координацию работы узлов.}
}

\newglossaryentry{id15}{
    name={Под (pod)},
    description={минимальная единица развертывания в Kubernetes, содержащая один или несколько контейнеров, совместно использующих ресурсы и сетевой контекст.}
}

\newglossaryentry{id16}{
    name={Признак (feature)},
    description={измеримая характеристика объекта или события, используемая при обучении и работе модели машинного обучения.}
}

\newglossaryentry{id17}{
    name={Пространство имён (namespace)},
    description={логическая область в Kubernetes, применяемая для изоляции и группировки ресурсов внутри кластера.}
}

\newglossaryentry{id18}{
    name={Сервисный аккаунт (service account)},
    description={учётная запись, используемая приложениями и компонентами системы для доступа к API и ресурсам Kubernetes.}
}

\newglossaryentry{id19}{
    name={Тестовая выборка},
    description={часть набора данных, не используемая при обучении модели и предназначенная для итоговой оценки качества её работы.}
}

\newglossaryentry{id20}{
    name={Функция потерь (loss function)},
    description={математическая функция, определяющая степень расхождения между предсказаниями модели и истинными значениями.}
}

\newglossaryentry{id21}{
    name={Эпоха обучения},
    description={один полный проход алгоритма обучения по всей обучающей выборке.}
}